% Pista geo-25j-q4b · Geometria
% Procedència: PAU Catalunya 2025 · Convocatòria de juny · Sèrie 1 · Exercici 4 — Opció B
\documentclass[11pt,a4paper]{article}
\usepackage{pista}
\pistainfo{Catalunya 2025}{Convocatòria de juny}{Sèrie 1}

\begin{document}

\section*{\textcolor{blauPista}{Exercici 4 --- Opció B}}

\noindent Considereu el pla $\pi$ d'equació $x + y = 0$.

\subsection*{\textit{a)} \normalfont Calculeu l'equació del pla $\pi'$ que és perpendicular a $\pi$ i conté els punts $P = (1, -1, 2)$ i $Q = (3, -3, 6)$. \hfill\textnormal{[1~punt]}}

\pista{Obtindràs l'equació general del pla a partir d'un punt i dos vectors directors. Pren el punt $P$; un dels vectors directors serà $\vec{PQ}$. L'altre és fàcil d'obtenir: com que els dos plans han de ser perpendiculars, el vector normal de $\pi$, $\vec{n}=(1, 1, 0)$, és un vector director de $\pi'$. Finalment, calcula el determinant que té per columnes $(x-1,\, y+1,\, z-2)$, $\vec{PQ}$ i $\vec{n}$. Recorda que et demanen una equació: no n'hi ha prou de calcular el determinant, has d'escriure $\det(\dots)=0$.}

\subsection*{\textit{b)} \normalfont Calculeu l'equació paramètrica de la recta continguda en $\pi'$ i que conté els punts de $\pi'$ a la mateixa distància de $P$ que de $Q$. \hfill\textnormal{[1,5~punts]}}

\pista{El conjunt de punts que estan a la mateixa distància de $P$ i $Q$ és un pla, que s'anomena el \emph{pla mediador} del segment $PQ$ (és un pla que passa pel punt mig entre $P$ i $Q$, i que és perpendicular a $\vec{PQ}$). La recta que et demanen és la intersecció d'aquest pla mediador amb $\pi'$: resol el sistema format per les dues equacions i expressa la solució en forma paramètrica.}

\end{document}
